\documentclass[10pt]{article}
\usepackage[margin=20mm]{geometry}
\usepackage{amsmath,amssymb,mathtools,fontspec,unicode-math,xurl,hyperref}
\IfFontExistsTF{DejaVu Serif}{\setmainfont{DejaVu Serif}}{\setmainfont{Latin Modern Roman}}
\setmathfont{latinmodern-math.otf}
\setlength{\parindent}{0pt}\setlength{\parskip}{6pt}\sloppy
\title{Initial Technical Review of
Regularity of Minimal Surfaces
with Capillary Boundary Conditions}
\author{AI-assisted technical review}
\date{}
\begin{document}
\maketitle
Abstract: The acquired manuscript is read through its definitions, main results, proofs, appendix and references. Five concrete mathematical corrections are established by elementary examples and direct differentiation: the model-selection implication in Lemma 3.15 admits an omitted wet-region branch; the variable-angle first-variation formula has the wrong derivative; the refined boundary flux uses the wrong trigonometric coefficient; two total-variation estimates omit a negative-mass contribution; and the boundary-normal linearization has an incorrect linear coefficient. The first issue is rated Major because an invalid model-selection step is explicitly used in the main regularity proof. The remaining issues have verified local corrections and are rated Minor. These findings do not establish failure of the principal regularity theorems. Scientific adjudication and independent acceptance remain pending.\par
\section*{Identity, version and scope}
The paper is by Luigi De Masi, Nick Edelen, Carlo Gasparetto and Chao Li, Regularity of minimal surfaces with capillary boundary conditions, DOI https://doi.org/10.1002/cpa.70008. This review concerns precisely arXiv:2405.20796v1, dated 31 May 2024, with 60 physical pages; physical and printed page numbers agree. The journal landing page identifies a 2025 publication, but the journal full-text request returned HTTP 403. Equivalence with the final journal text is not asserted. The current arXiv record and the author’s publication list did not supply a later arXiv revision or an identified correction in the bounded checks performed.\par
All 60 source pages were read as text. The displayed source pages supporting the findings were also visually inspected. Independent calculations concern the five examples below; reading the remaining arguments is not a claim to have independently reproved every external regularity or compactness theorem. No author code was executed. The scientific work and image inspection belong to this primary scientific conversation; native assistance was confined to mechanical binding and helper preparation.\par
\section*{Definitions and operative assumptions}
Equation reference: eq:capillary. Write \(W=\{x_1=0\}\), \(S_-=\{x\in W:x_{n+1}<0\}\) and \(S_+=\{x\in W:x_{n+1}>0\}\). The manuscript permits signed varifolds \[V=V_I+\cos\beta[S]_g,\qquad V_I\ \hbox{integral in }\{x_1<0\},\] where \(g,\beta\) are \(C^1\), \(S\) is measurable, and, for compactly supported \(C^1\) vector fields tangent to \(W\), \[
 \delta_{\beta,g}V(X)
 =\int\operatorname{div}_{V,g}X\,d\mu_{V_I}
  +\int_S\operatorname{div}_{W,g}(\cos\beta X)\,d\mathcal H_g^n
 =-\int g(H^{\rm tan},X)\,d\mu_{V_I}.\] This is Definition 1.1, p. 3, and Definition 2.1, pp. 13–14. In particular, negative wall weights are expressly allowed. The model is \[P_\theta^-=\{\cos\theta\,x_1+\sin\theta\,x_{n+1}=0,\ x_1\leq0\},
 \qquad V^{(\theta)}=[P_\theta^-]+\cos\theta[S_-].\] For \(0<\theta<\pi\), \[\Theta_{V^{(\theta)}}(0,1)+(\cos\theta)^-
   =\frac{1+|\cos\theta|}{2}<1.\] Here \(a^-=\max\{-a,0\}\). This shifted signed density must not be confused with normalized total variation.\par
The first-variation theorem requires a strict density gap \(1-\alpha\) and small \(H^{\rm tan}\), metric error and \(D\beta\). The main graph theorem excludes the right angle and requires the origin in the interior support and closeness to a half-plane. The version based on varifold distance also requires \(C^1\) closeness of the angle, not merely small \(D\beta\). The finite-perimeter conclusion for \(S\) is used only where \(|\cos\beta|\) stays away from zero. The dense regular-set assertion and the stronger almost-everywhere assertion in Theorem 1.19 have different density hypotheses. None of these restrictions is removed in this review.\par
\section*{Established problems}
Problem 1 (M25-0422-P001: The missing wet-region branch).  Severity: Major. Lemma 3.15, pp. 37–38, implication (3.38) and the angular conclusion, are false as stated. The omitted branch enters the first sentence of the proof of Theorem 1.11 on p. 54.\par
\subsection*{Claim and admissibility}
The lemma quantifies over \(\theta\) separated from \(0,\pi/2,\pi\), a \(C^1\) metric and angle, and a capillary varifold with \[0\in\operatorname{spt}V_I,\quad
 \Theta_V(0,1)+(\cos\beta(0))^-
 \leq\frac34+\frac{|\cos\theta_0|}{4},\] and arbitrarily small \(H^{\rm tan}\), metric error and \(D\beta\). It claims that small \(\operatorname{osc}_{P_\theta}(V_I,B_1)\) forces \(\mathcal D(V,V^{(\theta)})\) small and \(\beta\) close to \(\theta\), allowing addition of an integer multiple of \(\pi\).\par
Equation reference: eq:branch. Choose \(n=1\), \(g=g_{\rm eucl}\), \(\theta=\pi/3\), \(\theta_0=\pi/12\), \(\beta=2\pi/3\), and \[
 V=[P_{\pi/3}^-]-\tfrac12[S_+].\] Both angles are in their required ranges, \(V_I\) has multiplicity one, \(0\in\operatorname{spt}V_I\), and its oscillation relative to \(P_{\pi/3}\) is zero.\par
Exact counterexample. The outward conormal of the interior ray at the origin is \(\sin\theta\,e_1-\cos\theta\,e_2\). The conormal of the upper wall ray is \(-e_2\); multiplying by \(-\cos\theta\) cancels the tangential component. Thus [eq:capillary] (\#eq:capillary) holds with \(H^{\rm tan}=0\). All differential smallness quantities vanish. The normalized signed density is \(1/4\), so the shifted density is \(3/4\), strictly below \(3/4+\cos(\pi/12)/4\).\par
Nevertheless \[V-V^{(\pi/3)}=-\tfrac12[W]\ne0.\] Consequently their distance in the manuscript’s varifold metric is a fixed positive number. Taking a constant sequence of [eq:branch] (\#eq:branch) disproves (3.38) for any smaller prescribed distance tolerance and every positive smallness threshold. Moreover \[\inf_{k\in\mathbb Z}|2\pi/3+k\pi-\pi/3|=\pi/3,\] so the stated angular normalization does not rescue the conclusion. ◻\par
\subsection*{Dependency, counterevidence and repair boundary}
At the end of p. 37 the limiting wall coefficients satisfy \(k_\pm\in\{0,\cos\theta'\}\). Stationarity gives \(k_--k_+=\cos\theta\); it does not force the single solution asserted at the top of p. 38. Besides \((\cos\theta,0)\), the solution \((0,-\cos\theta)\) is allowed.\par
The manuscript itself supplies the relevant counterevidence to any claim of irreparable failure: Example 1.3 explicitly swaps the wet region without changing \(V_I\), and p. 40 announces an acute-angle normalization for Sections 4–6. The example is already a perfectly smooth half-plane, so it contradicts neither the graph conclusion nor capillary regularity. A corrected lemma should allow both signed models, or impose a compatible wet-region/sign normalization. For \(\beta,\theta\in(0,\pi)\), matching the signs of their cosines selects the intended branch; smallness of \(D\beta\) alone does not. The proof of Theorem 1.11 must make that selection before asserting the quoted convergence. This is a demonstrated invalid core reduction, not a demonstrated false final regularity theorem. A complete replacement proof with all parameter choices is not claimed here.\par
Problem 2 (M25-0422-P002: The derivative of the wall weight).  Severity: Minor. Equation (2.10), the definition of \(H_{\beta,g}\), and (2.12), p. 14, use \(\nabla_W^g\sin\beta\) where the product rule requires \(\nabla_W^g\cos\beta\). The second first-variation identity in Theorem 2.2, p. 15, is consequently false with the printed definition.\par
The hypotheses here are exactly the capillary definition, \(C^1\) metric and angle, integral interior part, and locally integrable tangential mean curvature. No density smallness is needed for this claim. In the Euclidean two-dimensional unit ball, take \(V_I=[P_{\pi/2}^-]\), \(S=W\), and \(\beta(y)=\pi/4+\kappa y\), where \(y=x_2\) and \(\kappa=1/10\). The wall weight is positive throughout the ball. The horizontal interior ray has zero tangential first variation, and the integral of \((\cos\beta X_2)'\) over the entire wall vanishes. Thus this is an admissible example with \(H^{\rm tan}=0\).\par
Equation reference: eq:weight. Opposite signs in the asserted identity. Let \(\phi(y)=(1-16y^2)^2\) for \(|y|<1/4\), extended by zero. This is a nonzero, nonnegative \(C^1\) function. Extend \(Y=(0,\phi(y))\) compactly into \(B_1\) with a normal cutoff equal to one near \(W\). The generalized boundary term is normal to \(W\) and hence vanishes on \(Y\). Directly, \[
 \delta V(Y)=\int\cos\beta\,\phi'\,dy
 =\kappa\int\sin\beta\,\phi\,dy>0.\] The printed \(H_\beta=(\nabla_W\sin\beta)/\cos\beta=\kappa e_2\) instead gives \(-\kappa\int\cos\beta\,\phi\,dy<0\). The two numerical values are approximately \(0.0188553390525\) and \(-0.0188553390525\); their signs follow analytically, without quadrature. ◻\par
The corrected identity is \[\delta_{\beta,g}V(X)=\delta_gV(X)
       +\int_S g(X,\nabla_W^g\cos\beta)\,d\mathcal H_g^n,\qquad
 H_{\beta,g}=\frac{\nabla_W^g\cos\beta}{\cos\beta}.\] Define the quotient to be zero on \(\{\cos\beta=0\}\) in the weighted almost-everywhere sense. The derivative vanishes almost everywhere on that level set, and \(|H_{\beta,g}|\,|\cos\beta|\leq|\nabla_W^g\beta|\). The required weighted \(L^1\) control therefore survives. The variable-angle compactness, monotonicity and first-variation estimates must use this corrected field; their small-error bounds are not disproved by the example. The analogous product-rule display in Remark 3.2, p. 25, should be corrected consistently. For constant \(\beta\), both derivatives vanish, explaining why constant-angle tests do not expose the error.\par
Problem 3 (M25-0422-P003: Normal flux in the refined boundary formula).  Severity: Minor. Theorem 3.16, p. 38, the text following (3.41), and the corresponding conclusion on p. 40 identify \(g(\eta_{V_I,g},\nu_g)\) with \(\cos\beta\). The coefficient is \(\sin\beta\) for the paper’s model.\par
Take \(n=1\), \(g=g_{\rm eucl}\), \(\beta=2\pi/3\), \(S=S_-\), \(V=V^{(\beta)}\), and \(\alpha=1/5\). The metric, curvature and angle-gradient errors are zero. The shifted density is \(3/4<4/5=1-\alpha\), so every smallness hypothesis of Theorem 3.16 holds. Its set \(BS\) contains the origin.\par
Model computation. For any \(0<\beta<\pi\), \[\eta_{V_I}=\sin\beta\,e_1-\cos\beta\,e_2,\qquad \nu=e_1.\] The lower wall ray contributes \(\cos\beta\,e_2\) at its endpoint. Thus \[\delta V(Y)=\sin\beta\,Y_1(0),\qquad
 \sigma_V=\sin\beta\,\delta_0 .\] At \(\beta=2\pi/3\) the actual positive coefficient is \(\sqrt3/2\), whereas the printed coefficient is \(-1/2\). The latter cannot be a nonnegative boundary measure. ◻\par
This is a false quantitative conclusion in the refined theorem, not an ambiguity about normal orientation: the outward normal is fixed to point in the positive \(e_1\) direction. For normalized angles in \((0,\pi)\), replace the coefficient on \(BS\) by \(\sin\beta\); the coefficient on its free-boundary complement remains one. This leaves the vector identity (3.41) and its geometric classification compatible with the model. Section 3.4 explicitly is not needed for the maximum principle or the main Allard-type regularity argument; the density classification used later in Theorem 1.19 is distinct from this flux coefficient. No failure of those qualitative conclusions follows from this example.\par
Problem 4 (M25-0422-P004: Total variation versus shifted signed mass).  Severity: Minor. Remark 3.3, (3.11), p. 25, and (3.33), p. 35, undercount the negative wall contribution when bounding \(|\mu_V|\).\par
Equation reference: eq:signedmass. Use \(n=1\), the Euclidean metric, the nonzero free-boundary horizontal ray \(V_I=[P_{\pi/2}^-]\), \(S=W\), and constant angle \(\cos\beta=-9/10\). This is capillary with \(H^{\rm tan}=0\); the full wall has no tangential boundary, and \(D\beta=0\). For every radius \(r\), with \(\omega_1=2\), \[
 \frac{\mu_V(B_r)}{\omega_1r}=-\frac25,\qquad
 \frac{|\mu_V|(B_r)}{\omega_1r}=\frac75,\qquad
 \Theta_V(0,r)+(\cos\beta)^-=\frac12.\] Taking \(\alpha=1/4\) verifies the strict density assumption \(1/2\leq3/4\), with zero differential errors. Nevertheless the conclusion (3.11) would bound \(7/5\) by \(7/8\), and (3.33) would bound \(7/5\) by one.\par
The exact measure identity is \[|\mu_V|=\mu_V+2\mu_V^- .\] For a wall density bounded below by \(-a\), a valid total-variation estimate uses \(\Theta_V+2a\), with the appropriate metric errors. In contrast, \(\mu_{V_I}\leq\mu_V+a\mathcal H_g^n\!\llcorner W\) does use only one copy of \(a\). The strict bound needed for interior boundary rigidity therefore survives. In Lemma 3.8, p. 33, the total mass controls an error of the form \(c\delta|\mu_V|(B_1)\); a uniform dimensional bound, rather than a bound strictly below \(\omega_n\), suffices for that use. Replacing the affected estimates by this valid bound repairs the identified mass accounting. This example does not refute the Ahlfors conclusions: the interior ray and its boundary are regular.\par
Problem 5 (M25-0422-P005: Incorrect linearization of the boundary normal).  Severity: Minor. The second Taylor expansion in Section 2.1, p. 10, has coefficient \(\cos^2\theta\sin\theta-1\) on \(D_1u\). It is reused in (5.7), p. 45, and in the boundary-viscosity argument of Theorem 6.3, p. 52.\par
Equation reference: eq:normal. The expansion concerns a \(C^2\) small slanted graph with fixed \(0<\theta<\pi\) and a small \(C^1\) metric perturbation. Set \(n=1\), \(g=g_{\rm eucl}\), \(\theta=\pi/3\), and \(u_t(x_1)=t x_1\). This satisfies all smallness requirements as \(t\to0\). By the upward unit-normal convention in Section 2.1, \[
 F(t):=\nu_\theta(u_t,g_{\rm eucl})\cdot e_1
 =\frac{\cot\theta-t}{\sqrt{1+(\cot\theta-t)^2}},\qquad
 F'(0)=-\sin^3\theta.\] Thus the residual after subtracting the printed linear term obeys \[\lim_{t\to0}
 \frac{F(t)-\cos\theta-(\cos^2\theta\sin\theta-1)t}{t}
 =1-\sin\theta>0.\] It cannot be absorbed in the claimed quadratic remainder. At \(\theta=\pi/3\) this limit is \(0.1339745962\ldots\). The independent difference quotients for \(t=10^{-2},\ldots,10^{-6}\) converge to it.\par
The correct first-order expansion is \(\cos\theta-\sin^3\theta\,D_1u+O_\theta(|Du|^2)\) in the Euclidean metric. Its coefficient remains strictly negative for a fixed angle away from the endpoints. Therefore the homogeneous Neumann condition obtained in the blow-up argument is unchanged; the barrier constants and displayed Taylor bounds must be adjusted. This is a disproved intermediate estimate with an explicit local repair, not a counterexample to the limiting elliptic problem.\par
\section*{Coverage, alternate readings and limitations}
The proof chain was followed from boundary first-variation control and rigidity (Sections 2–3 and Appendix 8), through the maximum principle (Section 4), partial Harnack estimates (Section 5), viscosity blow-up and decay (Section 6), to the main conclusions (Section 7). The role of the strict density gap, integral interior multiplicity, signed wall component, angle normalization, bounded mean curvature, metric regularity and non-right-angle restriction was checked at the corresponding transitions.\par
For Problem Problem 1 (\#prob:branch), merely changing the representative of \(\beta\) by a multiple of \(\pi\) cannot change the actual signed varifold. Swapping \(S\) and its complement can, and is the relevant repair. For Problems Problem 2 (\#prob:gradient)–Problem 5 (\#prob:normal), the counterexamples identify actual mathematical equalities or estimates that fail; their likely typographical origin does not make the printed formulas true. Their severities are limited because the stated corrections preserve the estimates or geometric conclusions at the specific uses examined.\par
The interior Allard statement on p. 23 also needs a clear centering convention if its \(C^{1,1/2}\) norm includes the graph height: the oscillation defined on p. 9 is invariant under normal translation. The notation for the zero generalized-boundary assumption there is ambiguous. This observation is not promoted to an additional unqualified counterexample or to a claim that Allard’s theorem fails. Uses at actual support points were read in their usual centered sense.\par
External results invoked for Allard regularity, rectifiability, compactness and the strong maximum principle were identified but not independently reproved from their original publications. No separate supplement resolving the findings was located in the bounded primary-source checks. The inaccessible journal full text remains a genuine version limit; the present findings must not be silently attributed to that unseen version. No claim of exhaustive absence of corrections is made.\par
\section*{Conclusion}
Review dimension | Recorded result\par
Acquired source | arXiv:2405.20796v1\par
Supported findings | One Major; four Minor\par
Principal regularity theorem disproved | No\par
Verified local repairs | Problems 2–5; branch correction identified\par
Complete replacement proof of Problem 1 | Not claimed\par
Author adoption of repairs | Unknown\par
Scientific adjudication | Pending\par
Independent acceptance | Not performed\par
\section*{Supporting checks}
The five proofs above are self-contained and are separately retained with the exact source locators, hypotheses and impact boundaries. The reviewer-authored calculation script used Python 3.13.5 and mpmath 1.3.0 at 70 decimal digits; it executed with exit code zero. Numerical quadrature and differentiation corroborate the analytic proofs and are not interval certificates.\par
The source PDF has 721,257 bytes and SHA256 09b84f9cabbf041c34753eb118aef10a2b448a0dc2f192c90e430a5a8bc9b382. The frozen manifest has SHA256 da5e45a8fc205eff01a855484924d74680139dae6d200a0b59fc496e3247b654. Exact report-source, compiler, dependency-payload, output, parse, render and later actual-view bindings are retained under `verification/0422`. Those records distinguish generation from actual inspection and warm-up from a final unchanged-dependency build.\par
9 L. De Masi, N. Edelen, C. Gasparetto and C. Li, Regularity of minimal surfaces with capillary boundary conditions, arXiv:2405.20796v1 (31 May 2024), 60 pp. https://arxiv.org/abs/2405.20796v1. Wiley, publication record for DOI 10.1002/cpa.70008, consulted 29 October 2026. https://onlinelibrary.wiley.com/doi/10.1002/cpa.70008. Metadata consulted; final full text unavailable in this review. N. Edelen, author publication list, consulted 29 October 2026. https://nedelen.science.nd.edu/.\par
\end{document}
